\logosection{\faGraduationCap}{教育经历}
\datedline{\textbf{福州大学} · 建筑材料与结构工程 · 硕士 · GPA \textbf{3.36/4.00}}{\dateRange{2024.09}{2027.06}}
\datedline{\textbf{湖南理工大学} · 土木工程 · 本科 · GPA \textbf{3.23/4.00}}{\dateRange{2020.09}{2024.06}}

\logosection{\faUniversity}{开源项目}
\datedline{\textbf{\href{https://github.com/Live-GalGame/LiveGalGame}{LiveGalGame} + \href{https://github.com/Live-GalGame/narrarc}{NarrowArc}} · 2.5k stars · 实时 AI 对话 Agent / 叙事型 RAG}{\dateRange{2025.12}{至今}}
\datedline{\biInfo{技术栈}{\textbf{TypeScript, Python, LangGraph, Qdrant, Reranker, RAG}}}{}
\begin{itemize}[itemsep=0.25ex, parsep=0ex]
    \item \textbf{实时决策控制:} 将流式 ASR 增量文本收敛为 committed turn，并以「轮到用户回应 / 连续多句 / 长时沉默」三类门控决定是否触发建议；通过双路音源分离维护双方上下文，避免 partial 文本误触发和角色混淆。
    \item \textbf{低延迟检索增强:} 采用文档 Embedding 异步预计算，在线链路按「Query Embedding $\to$ 本地 Qdrant Top 50 $\to$ Reranker Top 5 $\to$ LLM」执行；端到端 P90 控制在 \textbf{1.5s} 内，首 Token 延迟 \textbf{$<$ 900ms}。
    \item \textbf{跨会话长期记忆:} 构建「感知 $\to$ 反思 $\to$ 记忆写入 $\to$ 历史检索」闭环，持久化关系状态和关键事件；上线前后 30 天队列对比中，7 日留存由 \textbf{32\% 提升至 40\%}。
    \item \textbf{Agentic RAG 评测:} 为 NarrowArc 设计 Burst/TopicNode 双层索引及 planner-retriever-grader-explorer-generator 循环；在 RealTalk 50 条叙事问题上，闭环将 Phase Coverage 由 \textbf{55\% 提升至 85\%}，Fuzzy Recall 达 \textbf{78\%}，输出可回溯 evidence ids。
\end{itemize}

\datedline{\textbf{\href{https://www.date-match.online/}{DateMatch}} · AI 驱动社交匹配平台，非营利开源项目，用户规模近 20 万}{\dateRange{2026.01}{至今}}
\datedline{\biInfo{技术栈}{\textbf{PostgreSQL/pgvector, Kafka, Redis, Embedding, Reranker, LLM}}}{}
\begin{itemize}[itemsep=0.25ex, parsep=0ex]
    \item \textbf{分层决策工作流:} 负责「规则筛选 $\to$ Embedding 召回 $\to$ Rerank 相容性粗筛 $\to$ LLM 互补性决策」四阶段链路，将 20 万用户逐层收敛至 Top 200/Top 50；相较规则 + Embedding 基线，匹配接受率由 \textbf{18\% 提升至 24\%}。
    \item \textbf{向量召回优化:} 针对高选择性画像过滤场景选用 IVFFlat（lists=400, probes=30），结合标签预过滤与批量回表；20 万条 1024 维向量下 Recall@200 \textbf{$>$ 95\%}，召回链路 \textbf{P99 $<$ 80ms}。
    \item \textbf{可靠性与可观测性:} 通过 Kafka Topic/Consumer Group 解耦画像、Embedding、召回和决策，配置重试与死信队列；Redis 承担候选缓存、幂等和分布式锁，并以 Prometheus + trace\_id 跟踪分阶段延迟、错误率与消费积压。
\end{itemize}

\logosection{\faFlask}{科研经历}
\datedline{\textbf{基于 LAAF-BI-PINN 的输水隧道水锤安全风险评估}}{\dateRange{2026.01}{2027.06}}
\begin{itemize}[itemsep=0.25ex, parsep=0ex]
    \item 将控制方程、边界条件与观测数据统一纳入 PINN 损失，引入局部自适应激活函数和贝叶斯推断，完成关键响应预测与不确定性量化。
    \item 以 K-Means 划分运行状态，结合贝叶斯网络进行风险推理和敏感性分析，形成「响应预测 $\to$ 状态聚类 $\to$ 风险分级」闭环。
\end{itemize}

\logosection{\faCogs}{专业技能}
\begin{itemize}[parsep=0ex, itemsep=0ex, leftmargin=*]
    \item \textbf{Agent/LLM:} LangChain, LangGraph, RAG, Embedding/Rerank, Workflow Orchestration, Memory, Evals
    \item \textbf{工程能力:} Python, TypeScript, PostgreSQL/pgvector, Qdrant, Redis, Kafka, Prometheus, Git
\end{itemize}

\logosection{\faTrophy}{奖项荣誉}
\datedline{国家励志奖学金 · “高教杯”全国大学生先进成图大赛 · 校三好学生标兵 · 福州大学中期学业三等奖学金}{}
